\chapter{Shell e linha de comando}
\label{cap:shell}

\section{Vis\~ao geral}

O \code{src/shell.rs} implementa um terminal de linha de comando sem
aloca\c{c}\~ao din\^amica. Ele n\~ao \'e um shell POSIX, n\~ao interpreta
pipes, redirecionamentos, vari\'aveis, aspas ou scripts. \'E um la\c{c}o que:

\begin{enumerate}
  \item imprime o prompt \code{rinux> };
  \item l\^e bytes do teclado PS/2 ou da serial;
  \item ecoa caracteres imprim\'iveis e trata Backspace;
  \item ao receber Enter, executa o comando;
  \item volta ao in\'icio.
\end{enumerate}

\section{La\c{c}o principal}

\begin{lstlisting}[caption={Estrutura do la\c{c}o do shell.},label={lst:shellloop}]
pub fn run() -> ! {
    loop {
        let mut line = [0u8; LINE_MAX];
        let mut len = 0usize;
        let mut console = tty::active();
        crate::print!("rinux> ");

        loop {
            tty::pump_input();
            if tty::active() != console {
                console = tty::active();
                len = 0;
                crate::print!("rinux> ");
                continue;
            }

            let byte = match tty::read_byte() {
                Some(byte) => byte,
                None => { core::hint::spin_loop(); continue; }
            };

            match byte {
                b'\n' | b'\r' => {
                    println!();
                    execute(&line[..len]);
                    break;
                }
                0x08 | 0x7F => {
                    if len > 0 {
                        len -= 1;
                        crate::print!("\x08");
                    }
                }
                b'\t' => {}
                0x20..=0x7E => {
                    if len < LINE_MAX {
                        line[len] = byte;
                        len += 1;
                        crate::print!("{}", byte as char);
                    }
                }
                _ => {}
            }
        }
    }
}
\end{lstlisting}

\subsection{Buffer de linha}

A linha \'e um array fixo de 128 bytes na pilha do shell:

\begin{lstlisting}[caption={Buffer de linha.},label={lst:linebuffer}]
const LINE_MAX: usize = 128;
let mut line = [0u8; LINE_MAX];
let mut len = 0usize;
\end{lstlisting}

Quando a linha est\'a cheia, novos caracteres s\~ao silenciosamente
ignorados. N\~ao h\'a beep, mensagem de erro nem crescimento din\^amico.

\subsection{Entrada e eco}

A entrada \'e bombeada por \code{tty::pump\_input}, que drena o teclado PS/2 e
a serial COM1 para a fila do console virtual ativo. Ele tamb\'em interpreta
\code{Alt+F1}..\code{Alt+F8} como troca de console. O shell ent\~ao consome
bytes do console ativo com \code{tty::read\_byte}.

Se nenhum byte estiver dispon\'ivel, o shell executa \code{spin\_loop()} e
tenta novamente. N\~ao h\'a \emph{sleep}, j\'a que a espera \'e feita por
polling.

Caracteres imprim\'iveis (0x20 a 0x7E) s\~ao armazenados no buffer e ecoados
no console. Backspace (\code{0x08} ou \code{0x7F}) decrementa o tamanho e
imprime \code{\textbackslash x08}, que a camada TTY interpreta como apagar
a c\'elula anterior. Tab (\code{\textbackslash t}) \'e ignorado.

Se o console ativo mudar durante a edi\c{c}\~ao, a linha parcial \'e
descartada e o shell imprime um novo prompt no console selecionado.

\section{Interpreta\c{c}\~ao}

\subsection{Trim}

Antes de comparar, a linha passa por \code{trim}, que remove espa\c{c}os
ASCII no in\'icio e no fim. A fun\c{c}\~ao retorna uma \emph{slice} da
pr\'opria linha, sem c\'opia:

\begin{lstlisting}[caption={Remo\c{c}\~ao de espa\c{c}os nas extremidades.},label={lst:trim}]
fn trim(bytes: &[u8]) -> &[u8] {
    let start = bytes.iter()
        .position(|byte| !byte.is_ascii_whitespace())
        .unwrap_or(bytes.len());
    let end = bytes.iter()
        .rposition(|byte| !byte.is_ascii_whitespace())
        .map_or(start, |index| index + 1);
    &bytes[start..end]
}
\end{lstlisting}

Uma linha vazia ou composta apenas de espa\c{c}os n\~ao executa nada.

\subsection{Convers\~ao para texto}

\code{as\_text} tenta interpretar os bytes como UTF-8. Como o console s\'o
aceita ASCII, na pr\'atica a convers\~ao sempre funciona para texto digitado;
se falhar, retorna a string fixa \code{"<bytes invalidos>"}.

\subsection{Comandos}

A execu\c{c}\~ao \'e uma cadeia de compara\c{c}\~oes \code{if}/\code{else if}.
N\~ao h\'a tabela de despacho, \emph{parsing} de argumentos nem an\'alise
sint\'atica.

\begin{table}[H]
\centering
\caption{Comandos implementados.}
\label{tab:comandos}
\begin{tabularx}{\linewidth}{@{}lX@{}}
\toprule
\textbf{Comando} & \textbf{Efeito} \\
\midrule
\code{help} & Lista os comandos dispon\'iveis. \\
\code{clear} & Limpa o console virtual ativo com \code{ESC[2J ESC[H}. \\
\code{echo <texto>} & Imprime o texto ap\'os o primeiro espa\c{c}o. \\
\code{info} & Mostra informa\c{c}\~oes do kernel, TTY e modelo de entrada. \\
\code{tty} & Mostra o console virtual ativo. \\
\code{tty [0..7]} & Troca para o console virtual indicado. \\
\code{getty} & Roda o login em ring 3 no console ativo. \\
\code{getty [0..7]} & Troca de console e roda o login nele. \\
\code{demand} & Exercita o demand paging do kernel. \\
\code{usertest} & Roda o demo de demand paging em ring 3. \\
\code{userprot} & Tenta acessar o kernel a partir do ring 3. \\
\code{mem} & Mostra estat\'isticas de frames, faults e bits U/S. \\
\code{reboot} & Reinicia o PC via comando \code{0xFE} ao 8042. \\
\code{halt} & Imprime uma mensagem e para a CPU com \code{cli; hlt}. \\
\bottomrule
\end{tabularx}
\end{table}

\begin{nota}[Detalhes de comportamento]
\begin{itemize}
  \item \code{echo} sem argumento n\~ao \'e reconhecido como comando; apenas
        \code{echo } seguido de texto casa com o prefixo.
  \item \code{echo   ola} preserva os espa\c{c}os internos ap\'os o prefixo
        \code{"echo "}; o \emph{trim} remove apenas as extremidades.
  \item A compara\c{c}\~ao \'e sens\'ivel a mai\'usculas: \code{HELP} n\~ao
        funciona.
  \item Comandos desconhecidos produzem duas linhas de mensagem.
\end{itemize}
\end{nota}

\section{Sa\'ida de \code{info}}

O comando \code{info} imprime:

\begin{lstlisting}[style=shell,caption={Sa\'ida do comando \code{info}.},label={lst:info}]
Kernel: Rust, no_std, x86_64
Boot:   Multiboot2 (BIOS) -> modo longo
Video:  VGA texto 80x25, buffer 0xB8000
TTY:    8 consoles virtuais, ativo = tty0 (Alt+F1..Alt+F8)
Entrada: teclado PS/2 e serial COM1 (polling)
IDT:    256 portas; excecoes + IRQ0 (timer ~100 Hz)
Memoria: paginas de 4 KiB, demand paging e frames fisicos
Usuario: ring 3 com GDT/TSS, protecao U/S e int 0x80
Ticks do timer: 366
\end{lstlisting}

A maioria das linhas \'e est\'atica; o n\'umero do console ativo e a contagem
de ticks do timer s\~ao lidos no momento em que o comando roda.

\section{Comandos especiais}

\subsection{\code{halt}}

Imprime \code{"CPU parada. Use reset ou feche o QEMU."} e chama
\code{arch::halt\_loop()}, que executa \code{cli; hlt} indefinidamente. A
\'unica sa\'ida \'e reset f\'isico ou fechar o emulador.

\subsection{\code{reboot}}

Imprime \code{"Reiniciando..."} e chama \code{arch::reboot()}. O reset pode
n\~ao funcionar em todos os ambientes; se o comando ao 8042 for ignorado, a
CPU termina em \code{halt\_loop()}.

\section{Como estender}

Para adicionar um comando:

\begin{enumerate}
  \item acrescente um ramo \code{else if line == b"meucomando"} em
        \code{execute};
  \item use \code{println!} para a sa\'ida;
  \item se precisar de argumentos, use \code{strip\_prefix} como em
        \code{echo};
  \item mantenha o comando sem aloca\c{c}\~ao e sem bloqueio;
  \item atualize a lista de \code{help} e esta documenta\c{c}\~ao.
\end{enumerate}

Exemplo:

\begin{lstlisting}[caption={Esqueleto de um novo comando.},label={lst:novocomando}]
} else if line == b"versao" {
    println!("Rust VGA Kernel 0.1.0");
} else if let Some(arg) = line.strip_prefix(b"upper ") {
    for byte in arg {
        crate::print!("{}", byte.to_ascii_uppercase() as char);
    }
    println!();
}
\end{lstlisting}

\section{Limita\c{c}\~oes}

\begin{itemize}
  \item Sem hist\'orico, setas, Ctrl+C, Ctrl+D ou edi\c{c}\~ao no meio da
        linha.
  \item Sem aspas, escape, pipes, redirecionamento ou vari\'aveis.
  \item Sem argumentos estruturados: cada comando decide como interpretar o
        restante.
  \item Sem execu\c{c}\~ao de programas: n\~ao h\'a \emph{loader} nem
        processos.
  \item Sem tradu\c{c}\~ao de UTF-8: acentos digitados podem virar \code{?}.
  \item A linha \'e limitada a 128 bytes e o excesso \'e descartado em
        sil\^encio.
\end{itemize}
